% =====================================================================
% Appendix: Query Templates and Prompts
% =====================================================================
% Requires (in the main preamble):
%   \usepackage{tcolorbox}
%   \tcbuselibrary{breakable}
%   \usepackage{xcolor}
% Optional but recommended:
%   \usepackage{lmodern}     % nicer typewriter font
% =====================================================================

% -- Box style ---------------------------------------------------------
% Two visually-distinct boxes: one for raw LLM prompts (monospace), and
% one for natural-language example queries (regular text). Both are
% breakable so long prompts can flow across page boundaries.
\newtcolorbox{promptbox}[1]{%
    breakable,
    enhanced,
    colback=gray!4,
    colframe=gray!55!black,
    boxrule=0.45pt,
    arc=2.5pt,
    left=6pt, right=6pt, top=5pt, bottom=5pt,
    title=\textsf{\small\bfseries #1},
    fonttitle=\sffamily\small\bfseries,
    coltitle=white,
    colbacktitle=gray!55!black,
    attach boxed title to top left={yshift=-2mm, xshift=4mm},
    boxed title style={colback=gray!55!black, sharp corners, boxrule=0pt},
    fontupper=\small\ttfamily,
    before upper={\setlength{\parindent}{0pt}\setlength{\parskip}{2pt}},
}

\newtcolorbox{examplebox}[1]{%
    breakable,
    enhanced,
    colback=blue!2,
    colframe=blue!35!black,
    boxrule=0.45pt,
    arc=2.5pt,
    left=6pt, right=6pt, top=5pt, bottom=5pt,
    title=\textsf{\small\bfseries #1},
    fonttitle=\sffamily\small\bfseries,
    coltitle=white,
    colbacktitle=blue!35!black,
    attach boxed title to top left={yshift=-2mm, xshift=4mm},
    boxed title style={colback=blue!35!black, sharp corners, boxrule=0pt},
    fontupper=\small,
    before upper={\setlength{\parindent}{0pt}\setlength{\parskip}{2pt}},
}
% =====================================================================


\section{Query Templates and Prompts}
\label{app:prompts}

This appendix documents the textual prompts used throughout our benchmark.
We organise them into three parts: (i)~the prompts we feed to a large
language model to generate the query template banks once, before any
experiment is run; (ii)~representative examples of the resulting
\emph{standard} and \emph{explicit} queries on a subset of domains; and
(iii)~the prompts that each agent (or the shared domain classifier) issues
to a language model at inference time during the online evaluation.
Together these completely specify the linguistic stimuli that drive the
benchmark and the language-model behaviour we expect from each agent.


% =====================================================================
\subsection{Query Template Generation}
\label{app:prompts:generation}
% =====================================================================

We construct the benchmark's query bank in a single offline pass using a
strong language model (\textsc{GPT-5.2}). For each of the ten
\textsc{ToolBench-60} domains, we issue two distinct prompts. The
\emph{standard} prompt produces queries that describe a user's
information need in natural language but \emph{deliberately avoid
naming any specific tool}; such queries can be answered correctly only
if the agent has inferred the user's underlying preference. The
\emph{explicit} prompt produces queries that explicitly name one
particular tool from the domain's six options, so that a competent
language model can recover the target tool from the text alone. The
explicit queries serve as the \emph{out-of-distribution (OOD)} probes
used to evaluate override behaviour. We generate 20 standard templates
per domain and 1--5 explicit templates per tool; the exact counts and
the resulting raw text are released alongside the code.

\begin{promptbox}{Standard query generation prompt}
Generate \{n\_per\_domain\} diverse user queries that someone might ask
when needing a tool/API from the ``\{domain\}'' category.

Tools in this category:\\
\hspace*{1em}- \{tool\_1\}: \{description\_1\}\\
\hspace*{1em}- \{tool\_2\}: \{description\_2\}\\
\hspace*{1em}\ldots

IMPORTANT RULES:\\
\hspace*{1em}- Do NOT mention any specific tool name in the queries\\
\hspace*{1em}- Each query should be a natural user request (1-2 sentences)\\
\hspace*{1em}- Cover diverse use cases within this category\\
\hspace*{1em}- Write in English

Return a numbered list (1. query, 2. query, \ldots).
\end{promptbox}

\begin{promptbox}{Explicit query generation prompt}
Generate \{n\_per\_tool\} user queries that explicitly request using
``\{tool\}'' (a \{domain\} tool: \{description\}).

IMPORTANT RULES:\\
\hspace*{1em}- Each query MUST mention ``\{tool\}'' by name\\
\hspace*{1em}- The query should be a natural user request (1-2 sentences)\\
\hspace*{1em}- Write in English

Return a numbered list.
\end{promptbox}

The dual constraint---``do not mention any tool'' for standard queries
versus ``must mention this tool by name'' for explicit queries---creates
a clean binary partition of the benchmark's stimuli along the axis our
method is designed to handle: standard queries are decidable only via
\emph{personalised} statistical inference, while explicit queries are
decidable from the text alone. The standard-versus-explicit dichotomy
is the structural property that the two-bar test-accuracy decomposition
in the main paper makes visible.


% =====================================================================
\subsection{Example Queries by Domain}
\label{app:prompts:examples}
% =====================================================================

To give the reader a concrete sense of the linguistic style and
difficulty of the synthetic stimuli, we list three illustrative
domains from \textsc{ToolBench-60} below. For each domain we show three
standard queries (which require user-preference inference) and three
explicit queries that each name a distinct tool (which require only
text-level recognition). The full banks of 200 standard and 60
explicit templates across all ten domains are released with the
benchmark configuration.

\begin{examplebox}{Finance domain (6 tools)}
\textbf{Standard queries:}
\begin{itemize}\setlength{\itemsep}{1pt}
    \item Can you verify whether this routing and account number is
    valid and in good standing before I let the user submit an ACH
    payment? I want to reduce returns and fraud.
    \item I need to screen a batch of checking accounts before printing
    and mailing checks to new payees; can you flag risky accounts based
    on prior negative history?
    \item Pull real-time and 10-year historical daily prices for a list
    of US stocks, and also include splits and dividends where available.
\end{itemize}

\smallskip
\textbf{Explicit queries:}
\begin{itemize}\setlength{\itemsep}{1pt}
    \item Can you run an \emph{``Extended ACH and Check Prescreen''} on
    this routing and account number to confirm the account is open and
    in good standing before we accept an ACH payment?
    \item Use \emph{Twelve Data} to pull real-time and 1-minute intraday
    prices for AAPL for today, and return the latest quote plus OHLCV
    data.
    \item I want to integrate the \emph{ChangeNOW crypto exchange} into
    my checkout so customers can buy and exchange crypto; can you
    outline the steps and what I need to get started?
\end{itemize}
\end{examplebox}

\begin{examplebox}{Travel domain (6 tools)}
\textbf{Standard queries:}
\begin{itemize}\setlength{\itemsep}{1pt}
    \item I need to search and compare round-trip flight options from
    NYC to Tokyo for flexible dates next month, including price trends
    and the cheapest day to fly.
    \item Can you find the lowest fares for a multi-city itinerary
    (London $\to$ Rome $\to$ Athens $\to$ London) and return results
    with cabin class, baggage rules, and cancellation terms?
    \item I want to monitor a specific route (SFO to LAS) and alert me
    when the price drops below a threshold within the next 60 days.
\end{itemize}

\smallskip
\textbf{Explicit queries:}
\begin{itemize}\setlength{\itemsep}{1pt}
    \item Use the \emph{Skyscanner API} to find the cheapest round-trip
    flights from New York (JFK) to London (LHR) for flexible dates in
    the next two months, and return the top 10 options.
    \item Use \emph{Airbnb listings} to pull all active listings in the
    Lisbon market, including nightly price, cleaning fee, minimum stay,
    and availability for the next 60 days.
    \item Use \emph{webcams.travel} to find nearby landscape webcams
    around my current location, and show the top 10 with a quick
    preview and distance for each.
\end{itemize}
\end{examplebox}

\begin{examplebox}{Communication domain (6 tools)}
\textbf{Standard queries:}
\begin{itemize}\setlength{\itemsep}{1pt}
    \item I need to send and receive chat messages programmatically
    from a server, including handling delivery/read status callbacks.
    How can I integrate that into my app?
    \item Can you help me set up phone-number verification via SMS with
    a single API call, and automatically expire codes after five
    minutes?
    \item I have a user-entered phone-number string; I need it
    normalised into E.164 and also returned in national format with
    country details. What API can do that reliably?
\end{itemize}

\smallskip
\textbf{Explicit queries:}
\begin{itemize}\setlength{\itemsep}{1pt}
    \item Can you help me set up the \emph{WhatsApp Private API} to send
    and receive messages programmatically from my existing WhatsApp
    number?
    \item Please integrate \emph{``2Factor Authentication - India''}
    into my app for phone-number verification via OTP, using their
    single-API-call flow. Keep it fast and economical.
    \item Use \emph{Phone Formatter} to parse \texttt{+14155552671} and
    return the international and national formats, along with the
    detected country details.
\end{itemize}
\end{examplebox}

A few features are worth pointing out. Standard queries are
\emph{intentionally tool-agnostic}: a competent agent that has not
learnt the user's preference cannot do meaningfully better than uniform
random among the six in-domain tools. Explicit queries, by contrast,
quote the target tool name (in our shown examples we render the quoted
phrase in italics for readability), so a language model that can read
the query can identify the target tool without any history. The two
query types together cover the two sub-problems---personalisation and
explicit override---that motivate our architectural split.


% =====================================================================
\subsection{Inference-Time Prompts}
\label{app:prompts:inference}
% =====================================================================

At inference time, every agent that involves an LLM call issues one or
two structured prompts per round, each constrained to return a strict
JSON object. We list below the three prompts that drive the comparisons
reported in the main paper: the shared domain classifier (used by
\emph{all} agents that require an inferred domain, including
\textsc{Pure-Bandit}), the binary override probe used by
\textsc{Bandit-as-Override} and \textsc{Freq-as-Override}, and the
chain-of-thought selection prompt used by \textsc{Bandit-as-Context}.
For brevity we omit the prompts of the LLM-only baselines
(\textsc{ZeroShot-LLM}, \textsc{InContext-Memory},
\textsc{Profile-Memory}), which differ only in what user-history
context they prepend to a similar selection instruction.

\paragraph{Shared domain classifier.}
This prompt is issued exactly once per query and its inferred-domain
output is broadcast to every domain-aware agent. The placeholder
\texttt{\{domains\_block\}} expands to a multi-line listing of the ten
domains together with the first four tool names of each. Using a
shared classifier guarantees that all agents---statistical, LLM, and
hybrid alike---operate over the same domain signal, removing an
otherwise-unfair advantage that statistical agents would receive from
ground-truth domain access.

\begin{promptbox}{Shared domain classifier}
Which domain does this query belong to?

Domains and their tools:\\
\{domains\_block\}

Query: ``\{query\}''\\
Reply with only JSON: \{``domain'': ``<exact domain name>''\}
\end{promptbox}

\paragraph{Override probe.}
The defining design choice of our proposed \textsc{Bandit-as-Override}
agent (and its ablation \textsc{Freq-as-Override}) is to ask the LLM a
\emph{single binary question}: does the query explicitly name a tool?
The probe deliberately omits tool descriptions and includes only the
tool \emph{names} for the agent's inferred domain, so the LLM's task
reduces to lexical recognition rather than open-ended selection.
This narrow interface is what allows the LLM to act as an exception
handler without contaminating the bandit's default selection on the
$\sim$90\% of standard queries.

\begin{promptbox}{Override probe (Bandit-as-Override and Freq-as-Override)}
Does this query explicitly request a specific tool by name?\\
Available tools: \{tool\_1\}, \{tool\_2\}, \ldots, \{tool\_K\}\\
Query: ``\{query\}''

If the query mentions a specific tool name, reply:\\
\hspace*{1em}\{``override'': true, ``tool'': ``<tool name>''\}\\
If not, reply:\\
\hspace*{1em}\{``override'': false\}
\end{promptbox}

\paragraph{Selection-with-prior prompt.}
For comparison, the \textsc{Bandit-as-Context} baseline exposes the
bandit's posterior to the LLM as a tempered-softmax prior over the
domain's six tools, and asks the LLM to perform the \emph{full
selection} itself rather than answering an override question. The
placeholder \texttt{\{tools\_block\}} expands to lines of the form
``\texttt{- \{tool\_name\}: 38\%}'' sorted by descending prior. This is
the canonical ``LLM-as-Bayesian-reasoner'' design that our results show
to be both weaker and considerably more sensitive to the backbone's
reasoning quality.

\begin{promptbox}{Selection-with-prior prompt (Bandit-as-Context)}
Select the best tool for this query.\\
Available tools (with learned user preference \%):\\
\{tools\_block\}

Query: ``\{query\}''

If the query explicitly names a tool, use it. Otherwise, prefer the
tool with the highest user preference.\\
Reply with only JSON: \{``tool'': ``<tool name>''\}
\end{promptbox}

In all three prompts we set the decoding temperature to $0.0$ and parse
the response as JSON. On parse failure (empty body, malformed JSON, or
a tool name that does not appear in the domain's tool list), the
client first attempts a regular-expression and substring match against
the candidate tools, and only on total failure falls back to the
default action (the first listed domain for the classifier, the
bandit's selection for the override probe, and the first tool for
\textsc{Bandit-as-Context}). Empirically the fallback path fires on
$\lesssim 0.5\%$ of calls with the backbones reported in
\S\ref{ssec:hardware}.
